% Theoretical background: only the concepts the reader needs to follow the
% methodology and the results.

\chapter{Fundamentação teórica}
\label{chap:theory}

\placeholder{Apresentar, em um parágrafo, os temas abordados neste capítulo.}

\section{Sistemas embarcados no ensino}
\label{sec:embedded-education}

% Suggested sources: duarte2024, souza2014, souza2016, parreira2021, fruett2024.
\placeholder{Descrever o uso de plataformas embarcadas de baixo custo em atividades práticas de ensino e apresentar brevemente a placa usada como alvo de validação.}

\section{Modelos de linguagem de grande porte}
\label{sec:llm}

\placeholder{Conceituar LLM, seu uso na geração de código, o papel do contexto nos \textit{prompts} e a execução local de modelos abertos com ferramentas como o Ollama.}

\section{Aplicações \textit{web} e APIs REST}
\label{sec:web-rest}

\placeholder{Explicar a arquitetura cliente-servidor e o estilo arquitetural REST adotado na API.}

\section{Trabalhos relacionados}
\label{sec:related-work}

\placeholder{Comparar os trabalhos da base de referências com a proposta, evidenciando a lacuna tratada.}
